\newpage
\pagenumbering{arabic}
\setcounter{page}{1}
\markboth{}{} 

%\section*{\colorbox{lime}{Profesor}}
% \section*{Clase 14 de agosto}
% \addcontentsline{toc}{section}{Clase 14 de agosto}

\section*{Preliminares (de Topología)}

\begin{tcolorbox}[colback=blue-violet!5!white,colframe=blue-violet!75!black,title=\textbf{Definición}.]
    Un espacio \(X\) es \textit{Hausdorff\index{espacio Hausdorff}} si \(\forall\).
\end{tcolorbox}

\begin{tcolorbox}[colback=blue-violet!5!white,colframe=blue-violet!75!black,title=\textbf{Definición}.]
    Dado un espacio \(X\), una \textit{cubierta abierta (cerrada)\index{cubierta abierta}} es una familia de conjuntos abiertos (cerrados) \(\{U_\lambda\}_{\lambda\in\Lambda}\) tal que 
\end{tcolorbox}

\begin{tcolorbox}[colback=blue-violet!5!white,colframe=blue-violet!75!black,title=\textbf{Definición}.]
    Dado un espacio \(X\), una \textit{cubierta abierta (cerrada)\index{cubierta abierta}} es una familia de conjuntos abiertos (cerrados) \(\{U_\lambda\}_{\lambda\in\Lambda}\) tal que 
\end{tcolorbox}

\begin{tcolorbox}[colback=blue!5!white,colframe=blue!75!blue,title=\textbf{Teorema}.]
    Si \(A\subseteq X\), con \(A\) compacto y \(X\) es \(T_2\), entonces \(A\) es cerrado.
\end{tcolorbox}